\section{Warnungen des Kompilers}
\label{sec:c:warnungen}

FIXME

\cppinput{c/06_warnungen.c}

\paragraph{Erklärung:}

% \begin{description}
%     \item[\texttt{int zahl = 42;}]\hfill\\
%         Das ist die \textbf{Deklaration}\protect\footnote{%
%         } und \textbf{Initialisierung}\protect\footnote{%
%         } einer Variablen.

%         \texttt{int} ist der Datentyp, in diesem Fall eine ganze Zahl (Integer).\\
%         \texttt{zahl} ist der Name der Variable.\\
%         \texttt{42} ist der Wert, der der Variable zugewiesen wird.

%         Es ist auch möglich, eine Variable folgendermaßen zu deklarieren:
% \begin{code}
%   int zahl; // Deklaration ohne Initialisierung
% \end{code}
%         und ihr zu einem späteren Zeitpunkt einen Wert zuzuweisen:
% \begin{code}
%   zahl = 42; // Zuweisung eines Wertes
% \end{code}
%     \item[\texttt{printf("Die Zahl: \%d$\backslash$n", zahl);}]\hfill\\
%         Das ist die \textbf{Ausgabe}.

%         \texttt{printf} steht für \textit{print formatted} und schreibt Text
%         auf den Bildschirm (genauer: in die Konsole).

%         Das \texttt{\%d} ist ein Platzhalter für eine ganze Zahl, die an dieser Stelle ausgegeben wird.\\
%         Das \texttt{$\backslash$n} ist ein \textbf{Zeilenumbruch}.
%         Es sorgt dafür, dass der Cursor nach der Ausgabe in die nächste Zeile springt.
% \end{description}

\paragraph{Merke:}

% \begin{itemize}
%     \item Ein C-Programm beginnt immer in der \texttt{main()}-Funktion.
%     \item \texttt{printf()} schreibt Text auf den Bildschirm.
%     \item \texttt{$\backslash$n} macht einen Zeilenumbruch.
%     \item \texttt{return 0;} sagt dem Computer: ``Alles OK.''
%     \item Jede Anweisung endet mit \texttt{;}.
% \end{itemize}
